Source : calculs de l'auteur (fiscalite/is/calculs\_note\_is.py, mécanique de fiscalite/is/simulation\_is.py).\par Note : mêmes conventions que la figure précédente ; bâtiment amorti sur 25 ans, usure de 3,1 \% par an. « Petite société sous 42 500 \euro{} » : bénéfice inférieur au plafond du taux réduit, donc taxé à 15 \% à la marge aujourd'hui, à 25 \% avec la proposition. Proposition en 2032, quand la déduction atteint 100 \% au-delà du plafond, avant la baisse finale du taux (0,6 à 1,2 point, qui ramènerait le bâtiment vers 21 \%). Pour une société qui amortit aujourd'hui en dégressif, l'équipement est taxé à 15 \% au lieu de 20,7 \%. Les stocks et le besoin en fonds de roulement des petites sociétés suivent le même sort que leurs bâtiments.
